%===========================================
\documentclass[12pt,twoside]{article}
%===========================================
%\usepackage[latin1]{inputenc}
\usepackage[utf8]{inputenc}
\usepackage[T1]{fontenc}
\usepackage[spanish]{babel}
\usepackage[a4paper]{geometry}
%\geometry{top=1.5cm, bottom=1.0cm, left=1.25cm, right=1.25cm}
\geometry{
			a4paper,
			top=1.5cm, 
			bottom=1.0cm, 
			left=1.5cm, 
			right=1.25cm,
			headheight=14pt,
			includehead,
			includefoot}

\usepackage{amsmath, amsthm, amsfonts}
\usepackage{hyperref}
\usepackage{fancyhdr}
\usepackage{titlesec}
\usepackage{listings}
\usepackage{graphicx,graphics}
\usepackage{multicol}
\usepackage{multirow}
\usepackage{color}
\usepackage{subfig}
\usepackage[figuresright]{rotating}
\usepackage{enumerate}
\usepackage{anysize} 
\usepackage{url}
\usepackage{imakeidx}
\usepackage{float}
\usepackage{graphicx}
\usepackage{chngcntr}
\usepackage{enumitem}
% --- Contador global para \'items ---
\newcounter{globalitem}
\setcounter{globalitem}{0}

% --- Entorno que contin\'ua la numeraci\'on autom\'aticamente ---
\newenvironment{globalenum}{
  \begin{enumerate}[label=\arabic*., leftmargin=*, itemsep=2pt, topsep=2pt]
    % arranca en el \'ultimo valor usado
    \setcounter{enumi}{\value{globalitem}}
}{
    % guarda d\'onde nos quedamos
    \setcounter{globalitem}{\value{enumi}}
  \end{enumerate}
}

%===========================================
\newtheorem{Criterio}{Criterio}
\newtheorem{Sup}{Supuesto}
\newtheorem{Note}{Nota}
\newtheorem{Ejem}{Ejemplo}
\newtheorem{Prop}{Proposici\'on}
\newtheorem{Def}{Definici\'on}
\newtheorem{Teo}{Teorema}
\newtheorem{Algthm}{Algoritmo}
\newtheorem{Sol}{Soluci\'on}
\newtheorem{Cor}{Corolario}
\newtheorem{Ejer}{Ejercicio}
\newtheorem{Quest}{Pregunta}
\counterwithout{Ejer}{section}
\newtheorem{Obs}{Observaci\'on}
\newcommand{\rea}{\mathbb{R}}
\newcommand{\Prob}{\mathbb{P}}
\newcommand{\Esp}{\mathbb{E}}
\newcommand{\Var}{\mathbb{V}}
\newtheorem{Propty}{Propiedades}
\newtheorem{lem}{Lema}
%===========================================
\pagestyle{fancy}
\fancyhf{} % limpia encabezados y pies previos

\title{Ejercicios de Probabilidad y Probabilidad Condicional}
\author{}
\date{}

\begin{document}

\maketitle
\tableofcontents
%===========================================
\section{Tarea: Ejercicios para examen}
%===========================================

\begin{Ejer} Verify the following identities.

\begin{itemize}
\item[(a)]  $A \setminus B = A \setminus (A \cap B) = A \cap B^{c}$

\item[(b)] $B = (B \cap A) \cup (B \cap A^{c})$

\item[(c)] $B \setminus A = B \cap A^{c}$

\item[(d)] $A \cup B = A \cup (B \cap A^{c})$

\end{itemize}
\end{Ejer}

\begin{Ejer} Let $S$ be a sample space.

\begin{itemize}

\item[(a)] Show that the collection  $\mathcal{B}=\{\emptyset,S\}$ is a sigma algebra.

\item[(b)] Let $\mathcal{B} = \{\text{all subsets of } S, \text{ including } S \text{ itself}\}.$ Show that $\mathcal{B}$ is a sigma algebra.

\item[(c)] Show that the intersection of two sigma algebras is a sigma algebra.

\end{itemize}
\end{Ejer}

\begin{Ejer} Prove each of the following statements. (Assume that any conditioning event has positive probability.)

\begin{itemize}
\item[(a)] If $P(B)=1$, then $P(A|B)=P(A)$ for any $A$.

\item[(b)] If $A \subset B$, then $P(B|A)=1$ and $P(A|B)=\frac{P(A)}{P(B)}.$

\item[(c)] If $A$ and $B$ are mutually exclusive, then $P(A|A \cup B)=\frac{P(A)}{P(A)+P(B)}.$

\item[(d)] $P(A \cap B \cap C)=P(A|B \cap C)P(B|C)P(C).$

\item[(e)] $P(A\cap B) - P(A)P(B)=P(A^c)P(B) - P(A^c\cap B).$

\item[(f)] $P(A\cup B\cup C)=P(A)+P(A^c\cap B)+P(A^c\cap B^c\cap C).$

\end{itemize}
\end{Ejer}


\begin{Ejer}(Coding)

When coded messages are sent, there are sometimes errors in transmission. In particular, Morse code uses ``dots'' and ``dashes,'' which are known to occur in the proportion $3:4$. This means that for any given symbol, $P(\text{dot sent}) = \frac{3}{7}$ and $P(\text{dash sent}) = \frac{4}{7}.$ Suppose there is interference on the transmission line, and with probability $\frac18$ a dot is mistakenly received as a dash, and vice versa. If we receive a dot, can we be sure that a dot was sent?

Using Bayes Rule we write

$P(\text{dot sent}|\text{dot received})=P(\text{dot received}|\text{dot sent})\frac{P(\text{dot sent})}{P(\text{dot received})}$. 

We know $P(\text{dot sent})=\frac37$ and $P(\text{dot received}|\text{dot sent})=\frac78$.

Also

$P(\text{dot received})=P(\text{dot received} \cap \text{dot sent})+P(\text{dot received} \cap \text{dash sent})$

$=P(\text{dot received}|\text{dot sent})P(\text{dot sent})+P(\text{dot received}|\text{dash sent})P(\text{dash sent})$

$=\frac78 \cdot \frac37+\frac18 \cdot \frac47=\frac{25}{56}$.

Thus, $P(\text{dot sent}|\text{dot received})=\frac{(7/8)(3/7)}{25/56}=\frac{21}{25}$.
\end{Ejer}

\begin{Ejer} For any three events $A$, $B$, and $C$ defined on a sample space $S$, verify the following properties 

\textbf{Commutativity:} $A \cup B = B \cup A$ and $A \cap B = B \cap A$

\textbf{Associativity:} $A \cup (B \cup C) = (A \cup B) \cup C$ and $A \cap (B \cap C) = (A \cap B) \cap C$

\textbf{Distributive Laws:} $A \cap (B \cup C) = (A \cap B) \cup (A \cap C)$ and $A \cup (B \cap C) = (A \cup B) \cap (A \cup C)$

\textbf{DeMorgan's Laws:} $(A \cup B)^c = A^c \cap B^c$ and $(A \cap B)^c = A^c \cup B^c$
\end{Ejer}


\begin{Ejer}
Solve the following exercises
\begin{itemize}

\item[a) ] If $P(A)=\tfrac12$ and $P(B)=\tfrac12$, can $A$ and $B$ be disjoint? Explain.

\item[b) ] Prove or disprove: If $P(A)=P(B)=p$, then $P(A\cup B)\le p^2$.

\item[c) ] Prove or disprove: If $P(A)=P(B)$ then $A=B$.

\item[d) ] Prove or disprove: If $P(A)=0$, then $P(A|B)=0$.

\item[e) ] Prove or disprove: If $P(A)=\alpha$ and $P(B)=\beta$, then $P(A\cap B)\ge 1-\alpha-\beta$.

\end{itemize}
\end{Ejer}

\begin{Ejer} If $A$ and $B$ are two independent events defined on a given probability space $(\Omega,\mathcal{A},P)$, then

\begin{itemize}

\item[a) ] $A$ and $B^{c}$ are independent,

\item[b) ] $A^{c}$ and $B$ are independent,

\item[c) ] $A^{c}$ and $B^{c}$ are independent.

\end{itemize}
\end{Ejer}

\begin{Ejer} \textbf{Properties of $P(\cdot|B)$:} Assume the probability space $(\Omega,\mathcal{A},P)$ is given and let $B\in\mathcal{A}$ satisfy $P(B)>0$.

\begin{itemize}
\item[a) ] $P(\emptyset|B)=0.$

\item[b) ] If $A_1,\dots,A_n$ are mutually exclusive events in $\mathcal{A}$, then $P(A_1\cup\cdots\cup A_n|B)=\sum_{i=1}^{n}P(A_i|B).$

\item[c) ] If $A\in\mathcal{A}$ then $P(A^c|B)=1-P(A|B).$

\item[d) ] If $A_1,A_2\in\mathcal{A}$, then $P(A_1|B)=P(A_1A_2|B)+P(A_1A_2^c|B).$

\item[e) ] $P(A_1\cup A_2|B)=P(A_1|B)+P(A_2|B)-P(A_1A_2|B).$

\item[f) ] If $A_1\subset A_2$, then $P(A_1|B)\le P(A_2|B).$

\item[g) ]  $P(A_1\cup A_2\cup\cdots\cup A_n|B)\leq\sum_{i=1}^{n}P(A_i|B).$
\end{itemize}
\end{Ejer}

\begin{Def} (Indicator Function:) Let $\Omega$ be any space with points $\omega$ and $A$ any subset of $\Omega$.   The indicator function of $A$, denoted $I_A(\cdot)$, is the function with domain $\Omega$ and codomain $\{0,1\}$ defined by

$I_A(\omega)=
\begin{cases}
1 & \text{if } \omega\in A \\
0 & \text{if } \omega\notin A
\end{cases}
$

Clearly $I_A(\cdot)$ indicates the set $A$.
\end{Def}

\begin{Ejer} Properties of Indicator Functions: Let $\mathcal{A}$ be a collection of subsets of $\Omega$.

\begin{itemize}
\item[i) ] $I_{A^c}(\omega) = 1 - I_A(\omega)$.
\item[ii) ]  $I_{A_1 \cap \cdots \cap A_n}(\omega)= I_{A_1}(\omega) I_{A_2}(\omega)\cdots I_{A_n}(\omega).$

\item[iii) ] $I_{A_1 \cup \cdots \cup A_n}(\omega)= \max \{ I_{A_1}(\omega), \dots, I_{A_n}(\omega) \}.$

\item[iv) ] $I_A^2(\omega) = I_A(\omega).$
\end{itemize}
\end{Ejer}


\begin{Ejer} Solve the following
\begin{itemize}
\item[a) ] Show that $P(A|B) + P(A'|B) = 1,$ provided $P(B) \neq 0$.

\item[b) ] Show that $P(A \cap B \cap C \cap D)= P(A)\,P(B|A)\,P(C|A \cap B)\,P(D|A \cap B \cap C),$ provided $P(A \cap B \cap C) \neq 0$.

\item[c) ] Show that if $A$ is independent of $B$ then $A$ is also independent of $B'$, namely, that $P(A|B)=P(A)$ implies $P(A|B')=P(A),$ provided $P(B)\neq1$. 

\textit{Hint:} Make use of the fact that
$A=(A\cap B)\cup(A\cap B')$ where $A\cap B$ and $A\cap B'$ are mutually exclusive.

\item[d) ] Demonstrate that $P(A\cap B\cap C)=P(A)\,P(B)\,P(C)$ does not necessarily imply that $A$, $B$, and $C$ are independent.

\end{itemize}
\end{Ejer}

\begin{Ejer} Show that if set $B$ is contained in set $A$ (that is, $B$ is a subset of $A$), then
$P(B)\le P(A).$

\end{Ejer}

\begin{Ejer}

\begin{itemize}
\item[a) ] If $A$ and $B$ are mutually exclusive events, $P(A)=0.33$, and $P(B)=0.42$, find

\begin{itemize}
\begin{multicols}{2}
\item[i) ] $P(A')$;
\item[ii) ] $P(B')$;
\item[iii) ] $P(A\cup B)$;
\item[iv) ] $P(A\cap B)$;
\item[v) ] $P(A\cap B')$;
\item[vi) ] $P(A'\cap B')$.
\end{multicols}
\end{itemize}

\item[b) ] Given $P(A)=0.48$, $P(B)=0.37$, and $P(A\cap B)=0.13$, find

\begin{itemize}
\begin{multicols}{2}
\item[i) ] $P(A\cup B)$;
\item[ii) ] $P(A\cap B')$;
\item[iii) ] $P(A'\cup B')$;
\item[iv) ] $P(A'\cap B')$.
\end{multicols}
\end{itemize}

\item[c) ] Show that

\begin{itemize}
\begin{multicols}{2}
\item[i) ] $P(A\cap B) \le P(A)+P(B)$;
\item[ii) ] $P(A\cap B) \ge P(A)+P(B)-1$.
\end{multicols}
\end{itemize}

\item[d) ] Show that

\begin{align*}
P(A\cup B\cup C\cup D) &= P(A)+P(B)+P(C)+P(D) \\
&\quad -P(A\cap B)-P(A\cap C)-P(A\cap D) \\
&\quad -P(B\cap C)-P(B\cap D)-P(C\cap D) \\
&\quad +P(A\cap B\cap C)+P(A\cap B\cap D) \\
&\quad +P(A\cap C\cap D)+P(B\cap C\cap D) \\
&\quad -P(A\cap B\cap C\cap D).
\end{align*}

\end{itemize}
\end{Ejer}

\newpage

%===========================================
\section{Combinatoria}
%===========================================


\subsection{Propiedades de los N\'umeros Combinatorios}

Entre muchas propiedades de los n\'umeros combinatorios
exhibiremos las siguientes.

\begin{enumerate}
\begin{multicols}{2}
\item $\binom{n}{k}=\binom{n}{n-k}$

\item $\binom{n+1}{k}=\frac{n+1}{n-(k-1)}\binom{n}{k}$

\item $\binom{n+1}{k+1}=\frac{n+1}{k+1}\binom{n}{k}$

\item $\binom{n}{k+1}=\frac{n-k}{k+1}\binom{n}{k}$

\item $\binom{n}{k}+\binom{n}{k+1}=\binom{n+1}{k+1}$
\end{multicols}
\end{enumerate}

\subsection{Teorema del Binomio}

\begin{Teo}
Si $n\in\mathbb{N}$ y $a,b\in\mathbb{R}$ entonces $(a+b)^n=\sum_{k=0}^{n}\binom{n}{k}a^{n-k}b^k$
\end{Teo}

\begin{Cor}
El t\'ermino de orden $k+1$ es $t_{k+1}=\binom{n}{k}a^{n-k}b^k$
\end{Cor}

\begin{Ejer}

\begin{itemize} Resolver los siguientes ejercicios:

\item[a) ] En el desarrollo binomial $\left(x-\frac{1}{x^3}\right)^n$ demostrar que si aparece un t\'ermino de la forma $x^{-4m}$, entonces $n$ es m\'ultiplo de $4$.

\item[b) ] Determinar el t\'ermino independiente de $(2x+1)\left(1+\frac{2}{x}\right)^n$

\item[c) ] Demostrar que $\sum_{k=0}^{n}(-1)^k\binom{n}{k}=0$

\item[d) ] Si $A=(x^2+\frac{1}{x})^n,\qquad B=(x^3+\frac{1}{x^2})^n$ y el t\'ermino $k$-\'esimo coincide, demostrar que $n$ es par.

\end{itemize}
\end{Ejer}

\begin{Ejer}
\begin{itemize}
\item[i) ] Determinar el s\'eptimo t\'ermino de $(2x-y)^{12}$

\item[ii) ] Determinar el noveno t\'ermino de $\left(2+\frac{x}{4}\right)^{15}$

\item[iii) ] Determinar el t\'ermino 14 de $\left(4x^2y-\frac{1}{2xy^2}\right)^{20}$

\item[iv) ] Determinar el t\'ermino que contiene $x^2$ en $\left(\sqrt[3]{x}-\frac{2}{x^2}\right)^{27}$

\item[v) ] Determinar el t\'ermino que contiene $\frac{x^2}{y^2}$ en $\left(\frac{x}{y}-\frac{y^2}{2x^2}\right)^8$

\item[vi) ] Determinar el t\'ermino que contiene $x^r$ en $\left(x+\frac{1}{x}\right)^n$

\item[vii) ] Si en $(2x^2-\frac{1}{x})^{60}$ existe un t\'ermino de la forma $ax^{-4}$,
determinar $a$.

\item[viii) ] Determinar el t\'ermino independiente de $(x^3-\frac{1}{x^2})^{30}$

\item[ix) ] Determinar $a$ en $\left(\frac{x}{a}+\frac{1}{x}\right)^{20}$ para que el t\'ermino independiente sea igual al coeficiente de $x^2$.

\item[x) ] En $(x\sqrt{x}+\frac{1}{x^2})^n$ el coeficiente binomial del tercer t\'ermino es mayor
que el del segundo en $44$ unidades.

\item[xi) ] Mostrar que el coeficiente del t\'ermino central de $(1+x)^{2n}$ es igual a la suma de los coeficientes de los dos t\'erminos centrales de $(1+x)^{2n-1}$.

\item[xii) ] Dados $(x^2+\frac{1}{x})^n,\quad(x^3+\frac{1}{x^2})^n$ determinar el conjunto de $n$ para los cuales los terceros t\'erminos coinciden.

\item[xiii) ] Si en $(1+x)^{43}$ los coeficientes en las posiciones $(2m+1)$ y $(m+2)$ son iguales, encontrar $m$.

\item[xiv) ] Determinar el coeficiente de $x^n$ en $(1-x+x^2)^n(1+x)^{2n+1}$

\item[xv) ] En $\frac{x}{a}-y^2)^{15}$ el t\'ermino que contiene $y^{22}$ tiene coeficiente $-\frac{455}{27}$ determinar $a$.

\item[ixv) ] Demostrar que $\sum_{i=0}^{n}\binom{n}{i}=2^n$

\item[xvii) ] Si $p+q=1$, demostrar que $r_k=\binom{n}{k}p^kq^{n-k}$ satisface $\sum_{k=0}^{n}kr_k=np$

\end{itemize}

\end{Ejer}

\subsection{Conteo}

\begin{Def}[Principio Fundamental de Conteo]
Si un experimento se puede llevar a cabo en $k$ etapas, de las cuales la primera se puede realizar en $n_1$ formas,  $n_2$ para la segunda etapa, y as\'i sucesivamente hasta $n_k$ para la $k$-\'esima etapa,  entonces el n\'umero total de formas en las que puede realizarse el experimento es $$n_{1}n_{2}\dots n_{4}$$.

\end{Def}


\begin{Def}[Permutaciones]
Para $n$ elementos distintos,  el total de formas 	en que se pueden ordenar es $P(n)=n!$.

\end{Def}


\begin{Teo}
\begin{eqnarray}
nPr=P\left(n,r\right)=n\left(n-1\right)\left(n-2\right)\cdots\left(n-r+1\right)=\frac{n!}{\left(n-r\right)!}
\end{eqnarray}
\end{Teo}

\begin{Teo}[Permutaciones con repetici\'on]
el n\'umero de permutaciones de $n$ objetos de los cuales $n_{1}$ son iguales, $n_{2}$ son iguales y $n_{k}$ son de la misma clase, es

\begin{eqnarray}
nPn_{1}n_{2}\cdot n_{k}=\frac{n!}{n_{1}!n_{2}!\cdots n_{k}!}
\end{eqnarray}

\end{Teo}


\begin{Def}[Ordenaciones con Repetici\'on]
Paraa $n$ elementos distintos, el n\'umero total de formas en que se pueden ordenar con repetici\'on arreglos de $k$ elementos es $n^{k}$.

\end{Def}

\begin{Def}[Ordenaciones sin Repetici\'on]
Para $n$ elementos distintos, el n\'umero total de posibles arreglos de $k$ elementos sin repetir es $$\frac{n!}{\left(n-k\right)!}.$$

\end{Def}


\begin{lem}
\begin{eqnarray}
\left(\begin{array}{c}
n\\
n-r
\end{array}\right)=\left(\begin{array}{c}
n\\
r
\end{array}\right)
\end{eqnarray}

\end{lem}

\begin{Teo}
\begin{eqnarray}
\left(\begin{array}{c}
n+1\\
r
\end{array}\right)=\left(\begin{array}{c}
n\\
r-1
\end{array}\right)+\left(\begin{array}{c}
n\\
r
\end{array}\right)
\end{eqnarray}
\end{Teo}


\begin{lem}
\begin{eqnarray}
\left(\begin{array}{c}
n\\
n_{1},n_{2},\ldots,n_{k}
\end{array}\right)=\frac{n!}{n_{1}!n_{2}!\cdots n_{k}!}
\end{eqnarray}

\end{lem}

\begin{Def}[Combinaciones]

\end{Def}


\begin{Teo}
\begin{eqnarray}
C\left(n,r\right)=\left(\begin{array}{c}
n\\
r
\end{array}\right)=\frac{nPr}{r!}=\frac{n!}{r!\left(n-r\right)!}
\end{eqnarray}
\end{Teo}

\begin{Ejer} Resolver los siguientes ejercicios de conjuntos y de elementos de conteo.
\begin{enumerate}

\item Para $A=\left\{1,2,3\right\}$, $A=\left\{2,4\right\}$ y $C=\left\{3,4,5\right\}$, hallar $A\times B\times C$, $A\times (B\cap C)$, $A\times (B\cup C)$.


\item Para $A=\left\{a,b\right\}$, $B=\left\{2,3\right\}$ y $C=\left\{3,4\right\}$, determinar $A\times (B\cup C)$, $(A\times B)\cup (A\times C)$, $(A\times B)\cap (A\times C)$ y $A\times (B)\cap C)$.

\item Determinar el conjunto potencia del conjunto $S=\left\{1,2,3\right\}$ as\'i como el n\'umero de elementos que tiene. Hacer lo mismo para $S=\left\{1,2,3,4,5\right\}$, y as\'i sucesivamente para $S=\left\{1,2,3,\ldots, N\right\}$.

\item Construir un diagrama de \'arbol para $A\times B\times C$, donde $A=\left\{1,2,3\right\}$,  $B=\left\{a,b,c\right\}$ y $C=\left\{a,c,e\right\}$.


\item Si $U=\left[0,1\right]$, $A=\left\{x:|x-\frac{1}{2}|<\frac{1}{4}\right\}$ y $B=\left\{x:|x-\frac{5}{6}|<\frac{1}{3}\right\}$, encuentre $A\cup B$, $A\cap B$, $A\cap B^{c}$, $B\cap A^{c}$, $\left(A\cap B\right)^{c}$ y $\left(A\cup B\right)^{c}$. Adem\'as encuentre el \'area de cada uno de los conjuntos anteriores.


\item Sea $A=\left\{\left(x,y\right):0\leq x\leq y\leq4\right\}$, $B=\left\{\left(x,y\right):0\leq y, 4-x\leq4\right\}$, encuentre el \'area de los siguientes conjuntos: $A\cap B$, $A\cup B$, $B\cap A^{c}$, $A\cap B^{c}$, $\left(A\cap B\right)^{c}$ y $\left(A\cup B\right)^{c}$.


\item Sea $A=\left\{\left(x,y\right): x^2+y^2\leq4\right\}$, $B=\left\{\left(x,y\right): x^2+y^2\geq1\right\}$, encuentre el \'area del conjunto $A\cap B$, $A\cup B$, $B\cap A^{c}$, $A\cap B^{c}$, $\left(A\cap B\right)^{c}$ y $\left(A\cup B\right)^{c}$.

\item Sea $A=\left\{\left(x,y\right): x^2-y^2\leq0\right\}$, $B=\left\{\left(x,y\right): x^2+y^2\geq2\right\}$,  $C=\left\{\left(x,y\right):x\in R, |y|\leq1\right\}$, encuentre el \'area de $A\cap B\cap C$.

\item Sea $A=\left\{x:x^4-16\leq0\right\}$, $B=\left\{x:x^3-x\leq0\right\}$. Encuentre $A\cap B$, $A\cup B$, $A-B$, $B-A$, $\left(A\cup B\right)-A$, $\left(A\cup B\right)-B$ y $\left(A\cup B\right)-\left(A\cap B\right)$, adem\'as de sus \'areas.

\item Sea $A=\left\{x:x\in R, y\in R, 1<x^{2}+y^{2}\leq2\right\}$, $B=\left\{(x,y):x\in R, y\in R, x\in\left[-3,0\right], y\in\left[-3,0\right]\right\}$, encuentre el \'area del conjunto $A\cap B$, $A\cup B$, $B\cap A^{c}$, $A\cap B^{c}$, $\left(A\cap B\right)^{c}$ y $\left(A\cup B\right)^{c}$
\end{enumerate}
\end{Ejer}

\begin{Ejer} Resolver los siguientes ejercicios

\begin{enumerate}
\item Si no se permiten repeticiones, dado el conjunto $A=\left\{1,2,3,4,5,6\right\}$, cu\'antos n\'umeros de 4 d\'igitos se pueden formar cu\'antos de estos son menores que $400$? cu\'antos son pares? cu\'antos son m\'ultiplos de 5?


\item De cuantas maneras se pueden acomodar 7 personas, si se quiere que:  3 de ellas siempre est\'en juntas?  3 de ellas nunca est\'en juntas?

\item De cu\'antas maneras se pueden acomodar 4 ni\~nas y 2 ni\~nos si se quiere qu se sienten todos los nin\~nos juntos y todas las ni\~nas juntas? se sienten de manera alternada ni\~nos y ni\~nas?


\item Cu\'antas permutaciones distintas pueden formarse con todas las letras de las siguientes palabras: estad\'istica, carrera,  murcielago, Probabilidad.


\item Desarrollar y simplificar:

\begin{itemize}
\begin{multicols}{2}
\item $\left(2x+y^2\right)^5$,
\item $\left(x^2-2y\right)^6$,
\item $\left(a+b+c\right)^4$,
\item $\left(a+b+c+d\right)^5$,
\end{multicols}
\end{itemize}

\item De c\'uantas maneras puede elegirse un comit\'e de 3 hombres y 5 mujeres, de un grupo de 7 hombres y 10 mujeres?

\item Un estudiante tiene que responder 8 de 10 preguntas en un examen, cu\'antas maneras de responder tiene si 3 son obligatorias? si debe de responde 4 de las 5 primeras preguntas?


\item Un estudiante debe de responder un examen de 10 preguntas de opci\'on m\'ultiple: $\left\{a,b,c\right\}$. Cu\'antas maneras tiene de responder si lo hace al azar?. 


\item De cuantas maneras pueden repartirse 20 estudiantes en 3 grupos, $A_{1},A_{2}$ y $A_{3}$,   si todos deben de tener el mismo n\'umero de estudiantes? si la \'unica condici\'on es que haya m\'as de dos estudiantes?

\item Simplifique las siguyientes expresiones:

\begin{itemize}
\begin{multicols}{4}
\item $\frac{(n+1)!}{n!}$,
\item $\frac{n!}{(n-2)!}$,
\item $\frac{(n-1)!}{(n+2)!}$,
\item $\frac{(n-r+1)!}{(n-r-1)!}$.
\end{multicols}
\end{itemize}

\end{enumerate}

\end{Ejer}

\begin{Ejer} Resolver los siguientes ejercicios:
\begin{enumerate}

\item Para $\Prob\left[A-B\right]=\Prob\left[B-A\right]$, $\Prob\left[A\cup B\right]=1/2$ y $\Prob\left[A\cap B\right]=1/4$, encontrar $\Prob\left[B\right]$ y $\Prob\left[A^{c}\cap B\right]$.

\item Para $\Prob\left[A^{c}\cap B^{c}\right]=1/2$, $\Prob\left[A^{c}\right]=2/3$ y $\Prob\left[A\cap B\right]=1/4$, hallar $\Prob\left[B\right]$ y $\Prob\left[A^{c}\cap B\right]$.

\item Para $\Prob\left[A\right]=1/4$, $\Prob\left[B\right]=3/4$ y $A\cap B=\emptyset$, calcular $\Prob\left[A\cup B\right]$, $\Prob\left[A^{c}\cup B\right]$ y $\Prob\left[A\cup B^{c}\right]$.

\item Para $\Prob\left[A\cup B\right]=\Prob\left[B\right]=1/2$, $\Prob\left[A\cap B\right]=1/4$. Calcular $\Prob\left[A\right]$, $\Prob\left[A^{c}\right]$ y $\Prob\left[B\right]$.

\item Si se lanzan dos dados,  cu\'{a}l es la probabilidad de que las caras que muestren:
\begin{enumerate}
\item sumen seis?
\item su producto sea seis?
\item el valor de la diferencia sea uno?
\end{enumerate}

\item De una urna, con tres bolas blancas, dos negras y cuatro verdes, se sacan tres sin reposici\'{o}n. Explique la probabilidad de obtener:
\begin{enumerate}
\item bolas de tres colores;
\item bolas de un color;
\item dos bolas blancas;
\item por lo menos una bola blanca.
\end{enumerate}

\item En una competencia de nataci\'{o}n intervienen tres j\'{o}venes que llamaremos $A,B$ y $C$. Por su trayectoria en este tipo de competencias, se sabe que la probabilidad de que $A$ gane es el doble de la de $B$, y la probabilidad de que $B$ gane es el triple de la de $C$. Calcular:
\begin{enumerate}
\item Las probabilidades $\Prob\left[A\right]$, $\Prob\left[B\right]$ y $\Prob\left[C\right]$ que tiene cada joven de ganar.
\item La probabilidad de que $A$ no gane.

Sugerencia: $\Prob\left[A\cup B\cup C\right]=\Prob\left[A\right]+\Prob\left[B\right]+\Prob\left[C\right]-\Prob\left[A\cap B\right]-\Prob\left[A\cap C\right]-\Prob\left[B\cap C\right]+\Prob\left[A\cap B\cap C\right]$
\end{enumerate}

\item Si una persona acude con su dentista, supongamos que la probabilidad de que le limpie la dentadura es de $0.44$, la probabilidad de que le tape una caries es de $0.24$, la probabilidad de que se le extraiga un diente es  $0.21$, la probabilidad de que se le limpie la dentatura y le tape una caries es de $0.08$, la probabilidad de que le limpie la dentadura y le extraiga un diente es $0.11$, la probabilidad de que le tape una caries y le saque un diente es de $0.07$, y la probabilidad de que le limpie la dentadura, le tape una caries y le saque un diente es $0.03$.  Cu\'{a}l es la probabilidad de que una persona que acudea su dentista se le haga por lo menos una de estas cosas?

\item Un equipo de baloncesto consta de 6 jugadores negros y 6 jugadores blancos. Para acomodarlos en sus habitaciones se van a elegir pares de jugadores. Si estos pares se forman al azar,  cu\'l es la probabilidad de que ninguno de los jugadores negros tengan un compa\~nero de cuarto blanco? Sugerencia: contar primero cuantas maneras posibles hay de encontrar pares solamente con jugadores negros, luego cuantas maneras distintas hay de encontrar pares de jugadores blancos y aplicar la regla fundamental del conteo, finalmente determinar cuantos posibles pares hay considerando elecciones sin reemplazo y aplicar nuevamente la regla fundamental del conteo.

\item En la urna I hay tres bolas blancas y 5 negras. En la II hay cinco blancas y tres negras. De cada urna se separa una bola. Determine la probabilidad de obtener:
\begin{enumerate}
\item dos bolas blancas;
\item bolas de diferentes colores,
\item por lo menos una bola blanca.
\end{enumerate}
\end{enumerate}
\end{Ejer}

\begin{Ejer} Resolver los siguientes ejercicios: 
\begin{enumerate}
\item Hay tres urnas: en la I se han puesto tres bolas blancas y dos negras; en la II hay cuatro bolas blancas y una negra; y en la III dos blancas y dos negras. De cada urna se escoge una bola. Determine la probabilidad de obtener:
\begin{enumerate}
\item tres bolas blancas;
\item una bola blanca,
\item por lo menos una bola blanca.
\end{enumerate}

\item En m\'exico la probabilidad, que un directivo sea egresado de una universidad privada o que tenga maestr\'ia, es de $60\%$. De que sea egresado de una universidad privada es $20\%$ y la de que tenga una maestr\'ia es de $40\%$,  cu\'al es la probabilidad de que uno escogido al azar sea egresado de una universidad privada y tenga maestr\'ia?

\item Los ahorradores de una peque\~na poblaci\'on cuentan con tres bancos: A, B y C. Suponga que $60\%$ de las familias de la ciudad han depositado sus ahorros en el banco $A$, $40\%$ en el B y $30\%$ en el C. Tambi\'en, que $20\%$  de los habitantes abri\'o cuentas en A y B, $10\%$  en A y C; $20\%$  en B y C y $5\%$ en A, B y C.  Qu\'e porcentaje de las familias de la ciudad manejan sus cuentas en al menos uno de los tres bancos?

\item Despu\'es de realizar un estudio en una universidad particular, en el \'area de posgrado, se estim\' que $30\%$  de los estudiantes estaban seriamente preocupados por sus posibilidades de encontrar trabajo, $25\%$  por sus notas y $20\%$ por ambas cosas.  cu\'al es la probabilidad de que un estudiante del \'ultimo curso, elegido al azar, est\' preocupado por al menos una de las dos cosas?

\item Si se tiran tres dados,  cu\'al es la probabilidad de que en dos aparezca el mismo n\'umero?

\item Si se lanzan tres dados y dos monedas,  cu\'antas probabilidades hay de obtener dos 6 y dos soles?

\item En un examen, un estudiante obtendr\'a una de las siguientes calificaciones: A, B, C y D. Con A, B y C, aprobar\'. La probabilidad de que saque A o B es de 0.6. Indique la probabilidad de que obtenga C si se sabe que la probabilidad de que apruebe el examen es de 0.9.

\item Se lanzan tres dados. Calcule la probabilidad de que la suma de los n\'umeros obtenidos sea mayor que 10, si la suma en los dos primeros dados es cinco.

\item Un profesor de estad\'istica de una universidad imparte clases a un grupo, que consiste de 10 estudiantes de contabilidad, 14 de administraci\'on, 9 de tecnolog\'a de alimentos y 3 de mercadotecnia. Por experiencia, sabe que la probailidad de que un estudiante de conta, 14 de contabilidad, 14 de adminsitraci\'on, 9 de tecnolog\'ia de alimentos y 3 de mercadotecnia. Por experiencia sabe que la probabilidad de que un estudiante de contabilidad le copie la tarea a un compa\'ero, en lugar de resolverla, es de 0.9, para la administraci\'on, tecnolog\'ia de alimentos y mercadotecnia, las cifras correspondientes son de 0.7, 0.4 y 0.8, respectivamente. si al calificar una tarea descubre que las soluciones han sido copiadas,  cu\'al es la probabilidad de que se trate de la de un alumno de contabilidad?

\item El departamento de ventas de una compa\~n\'ia farmac\'eutica public\'o los siguientes datos relativos a las ventas de cierto analg\'esico fabricado por ellos
\begin{table}[!ht]
\begin{center}
\begin{tabular}{|c|c|c|}
\hline
 Analg\'esico & $\%$ de Ventas & $\%$ del grupo vendido en dosis fuerte \\\hline
C\'psulas & 57 & 38 \\\hline
Tabletas & 43 & 31 \\\hline
\end{tabular}
\end{center}
\end{table}
Si un cliente adquiri\'o la dosis fuerte de este medicamento,  cu\'al es la probabilidad de que lo comprara en forma de c\'apsulas?

\end{enumerate}
\end{Ejer}


\subsection{Probabilidad Condicional}


\begin{Ejer} Resolver la siguiente lista de ejercicios relacionados con Probabilidad Condicional e Independencia de Eventos
\begin{enumerate}


\item Un recipiente contiene 5 transistores defectuosos (fallan inmediatamente al ponerse en uso), 10 parcialmente defectuosos (fallan despu\'s de unas horas en uso) y 25 transistores aceptables. Del recipiente se toma de manera aleatoria un transistor y se pone en funcionamiento. Si no falla inmediatamente,  cu\'al es la probabilidad de que sea aceptable?


\item La organizaci\'on en la ue trabaja Luis va a ofrecer, a los empleados que tienen por lo menos un hijo var\'on, una comida para padre e hijo. Se invita a cada uno de estos empleados a que asista con su hijo menor. Si se sabe Luis tiene s\o'lo dos hijos,  cu\a'l es la probabilidad condicional de que ambos sean ni\~nos puesto que est\'a invitado a la comida? Suponga que el espacio muestral est\'a dado por $S=\left\{\left(b,b\right),\left(b,g\right),\left(g,b\right),\left(g,g\right)\right\}$ y que todos estos resultados sean igualmente posibles. Sugerencia: Sea B el evento de que los dos son varones y A el evento que por lo menos uno de ellos es var\'on.

\item La se\~nora P\'erez piensa que hay un $30\%$ de posibilidad de que la empresa donde labora abra una sucursal en San Luis. Si lo hace, ella tiene un $60\%$ de seguridad de que ser\' nombrada directora de esta nueva oficina.  Con que probabilidad la se\~nora P\'erez ser\' la directora de una sucursal en San Luis?



\item Si $\Prob\left[A^{c}\right]=1/3$, $\Prob\left[A\cap B\right]=1/4$ y $\Prob\left[A\cup B\right]=2/3$, calcule $\Prob\left[B\right]$, $\Prob\left[A\cap B^{c}\right]$ y $\Prob\left[B-A\right]$.

\item Si $\Prob\left[A^{c}\right]=1/3$, $\Prob\left[A\cup B\right]=5/6$ y $\Prob\left[B^{c}\right]=1/2$, calcule $\Prob\left[A\cap B\right]$, $\Prob\left[A\cap B^{c}\right]$ y $\Prob\left[B-A\right]$.

\item La sangre se clasifica por factores en $Rh$ positivo y $Rh$ negativo y tambi\'n de acuerdo al tipo. Si la sangre contiene un ant\'igeno $A$, es de tipo $A$; si tiene un ant\'igeno $B$, es de tipo $B$; si tiene ambos ant\'igenos $A$ y $B$, es de tipo $AB$, si carece de ambos ant\'igenos es de tipo $O$. En una encuesta a 75 individuos, se encontr\' que 65 de ellos ten\'ian factor $Rh^{+}$, de los cuales 25 con ant\'igeno $A$, 30 con ant\'igeno $B$ y 10 con ambos ant\'igenos. De los 10 con factor $Rh^{-}$ 3 ten\'ian ant\'igeno $A$, 4 con ant\'igeno $B$ y 1 con ambos ant\'igenos. Si se selecciona al azar a una de las 75 personas, calcule la probabilidad de que tenga sangre con:

\begin{enumerate}
\item Factor $Rh^{+}$ y tipo $A$,
\item Factor $Rh^{+}$ y tipo $AB$,
\item Factor $Rh^{+}$ y tipo $O$,
\item Factor $Rh^{-}$ y tipo B,
\item Factor $Rh^{-}$ y tipo $O$,
\end{enumerate}


\item Se sabe que un grupo grande de personas $10\%$ toman desayuno caliente, $20\%$ toman almuerzo caliente y $25\%$ toma desayuno cliente o almuerzo caliente. Encuentra la probabilidad de que una persona elegida al azar del grupo
\begin{enumerate}
\item tome desayuno caliente y un almuerzo caliente
\item Tome un desayuno caliente, dado que la persona elegida tom\' un desayuno caliente
\end{enumerate}

\item Si los eventos $A$ y $B$ son independientes y $\Prob\left[A\right]=0.3$, $\Prob\left[B\right]=0.5$, encuentra $\Prob\left[A\cap B\right]$ y $\Prob\left[A\cup B\right]$.

\item Los eventos $A$ y $B$ son tales que $\Prob\left[A\right]=2/3$, $\Prob\left[A|B\right]=2/3$, $\Prob\left[B\right]=1/4$. Encuentra $\Prob\left[A|B\right]$, $\Prob\left[B|A\right]$ y $\Prob\left[A\cap B\right]$.

\item En un grupo de 100 personas, 40 tienen un gato, 25 tienen un perro y 15 tienen un perro y un gato. Encuentra la probabilidad de que una persona elegida al azar
\begin{enumerate}
\item tenga un perro o un gato,
\item tenga un perro o un gato, pero no ambos.
\item tenga un perro dado que tiene un gato,
\item no tenga gato dado que tenga un perro.
\end{enumerate}
\end{enumerate}
\end{Ejer}

\begin{Ejer} Resolver la siguiente lista de ejercicios relacionados con Probabilidad Condicional e Independencia de Eventos
\begin{enumerate}

\item Sean $A$ y $B$ eventos exhaustivos y se sabe que $\Prob\left[A|B\right]=1/4$ y $\Prob\left[B\right]=1/3$. Encuentra $\Prob\left[A\right]$.

\item Una bolsa contiene 4 canicas rojas y 6 negras. Se saca una canica al azar y no se regresa. Entonces se saca otra canica. Encontrar la probabilidad de que
\begin{enumerate}
\item la segunda canica sea roja, dado que la primera fue roja
\item las dos canicas sean rojas,
\item las canicas sean de diferente color.
\end{enumerate}

\item Sean $A$ y $B$ dos eventos independientes tales que $\Prob\left[A\right]=0.2$, $\Prob\left[B\right]=0.15$, Evalue las siguientes probabilidades
\begin{enumerate}
\item $\Prob\left[A|B\right]$
\item $\Prob\left[A\cap B\right]=1/4$
\item $\Prob\left[A\cup B\right]=1/4$.
\end{enumerate}

\item Sean $A$ y $B$ dos eventos tales que $\Prob\left[A\right]=0.4$, $\Prob\left[B\right]=0.25$. Si $A$ y $B$ son independientes evalue las siguientes probabilidades
\begin{enumerate}
\item $\Prob\left[A\cap B\right]$
\item $\Prob\left[A\cap B^{c}\right]$
\item $\Prob\left[A^{c}\cap B^{c}\right]$.
\end{enumerate}
	
\item Si $\Prob\left[A^{c}\right]=1/3$, $\Prob\left[B|A\right]=1/2$, calcule $\Prob\left[B\cap A\right]$.

\item Si $\Prob\left[A\right]=1/4$, $\Prob\left[B|A\right]=1/4$, calcule $\Prob\left[B^{c}\cap A\right]$.

\item Si $\Prob\left[A\right]=1/3$, $\Prob\left[A|B\right]=1/5$, y $\Prob\left[B|A\right]=1/2$, calcule $\Prob\left[B\right]$.

\item Si $\Prob\left[B\right]=\Prob\left[B^{c}\right]$, $\Prob\left[A|B^{c}\right]+\Prob\left[A|B\right]=4/5$, determine $\Prob\left[A\right]$.

\item Si A y B son independientes y $\Prob\left[A^{c}\right]=1/2$. $\Prob\left[A\cap B\right]=1/5$, determine $\Prob\left[B\right]$.

\item Si A y B son independientes y $\Prob\left[A\right]>0$, $\Prob\left[A\right]=3\Prob\left[A\cap B\right]$, determine $\Prob\left[B\right]$.

\end{enumerate}
\end{Ejer}

\begin{Ejer} Resolver la siguiente lista de ejercicios:
\begin{enumerate}
\item Si A, B y C son independientes y $\Prob\left[A\cap B\right]=1/2$, $\Prob\left[C^{c}\right]=1/3$, determine $\Prob\left[A\cap B\cap C\right]$.

\item Si A, B y C son independientes y $\Prob\left[A\cap B\cap C\right]=1/8$, $\Prob\left[A\right]=2\Prob\left[B\right]=4\Prob\left[C\right]$, determine $\Prob\left[C^{c}\right]$.

\item Si A y B son independientes y $\Prob\left[A\right]>0$, $2\Prob\left[A\right]=5\Prob\left[A\cap B\right]$, determine $\Prob\left[B^{c}\right]$.

\item Si A, B, C y D son independientes y $\Prob\left[A\cap B\cap C\right]=1/24$, $\Prob\left[A\right]=2\Prob\left[B\right]=3\Prob\left[C\right]=4\Prob\left[D\right]$, determine $\Prob\left[D\right]$.

\item \'Pueden ser mutuamente excluyentes los eventos A y B, si $\Prob\left[A\right]=0.54781$ y $\Prob\left[B\right]=0.49719$

\item Si los eventos A y B son independientes y $\Prob\left[A\right]=0.3$ y $\Prob\left[B\right]=0.5$, encuentra $\Prob\left[A\cap B\right]$ y $\Prob\left[A\cup B\right]$.

\item A y B son eventos exhaustivos y se sabe que $\Prob\left[A|B\right]=1/4$ y $\Prob\left[B\right]=2/3$. Encuentra $\Prob\left[A\right]$.

\item Sean A y B dos eventos independientes tales que $\Prob\left[A\right]=0.2$ y $\Prob\left[B\right]=0.15$. Eval\'ua las siguientes probabilidades: $\Prob\left[A|B\right]$, $\Prob\left[A\cap B\right]$ y $\Prob\left[A\cup B\right]$.

\item La probabilidad de que un evento A ocurra es $\Prob\left[A\right]=0.4$. B es un evento independiente de A y la probabilidad de la uni\'on de A y B es $\Prob\left[A\cup B\right]=0.7$. Encuentra $\Prob\left[B\right]$.

\item Los eventos A y B son independientes tales que $\Prob\left[A\right]=0.4$ y $\Prob\left[B\right]=0.25$. Determina: $\Prob\left[A\cap B\right]$, $\Prob\left[A\cap B^{c}\right]$ y $\Prob\left[A^{c}\cap B^{c}\right]$.

\end{enumerate}
\end{Ejer}


\subsection{Teorema de Probabilidad Total y Teorema de Bayes}

\begin{Ejer} Resolver los siguientes ejercicios

\begin{enumerate}

\item Una Compa\~n\'ia de Seguros divide a las personas en dos clases, quienes son propensos a accidentes y quienes no lo son. Sus estad\'isticas muestran que una persona propensa a accidentes, tendr\'a, en no m\'as de un a\~no, un accidente con probabilidad de $0.4$; mientras que esta probabilidad decrece a $0.2$ en personas no propensas a accidentes. Si pensamos que 30 por ciento de las personas es propenso a accidentes cu\'al es la probabilidad de que una persona que compra una nueva p\'oliza tenga un accidente en no m\'as de un a\~no?

\item Se tienen dos tarjetas, una de ellas es negra por ambas caras y la otra muestra una negra y otra blanca. Se meten en una bolsa y se extrae una de las dos al azar, la cual se coloca sobre la mesa. Si la cara hacia arriba es negra,  cu\'al es la probabilidad de que tambi\'en la de abajo sea negra?

\item Un rat\'on de laboratorio se introduce en un laberinto en forma de T. Del lado izquierdo, hay un pedazo de comida protegido para que el animal  no lo huela de lejos, en tanto que del lado derecho hay una peque\~na descarga el\'ectrica, que ser\'a desagradable para \'este, m\'as no mortal. Suponga que la primera vez que se mete el rat\'on hay una probabilidad de 0.5 de que vira a cualquiera de los dos lados. Si en el primer intento fue a la izquierda, entonces hay una probabildad de 0.6 de que vuelva a hacerlo en el segundo; sin embargo, si en el primero dio vuelta a la derecha, y recibi\'o la peque\~na descarga el\'ectrica, entonces hay una probabilidad de 0.75 de que se ir\'a la izquierda en el segundo. Si se observa que el animal efectivamente camin\'o a la izquierda en el segundo intento,  cu\'al es la probabilidad de que haya virado tambi\'en hacia el mismo lado en el primer intento?

\item En cierto pa\'is (no es M\'xico) aquejado por la inflaci\'on, los economistas proponen tres teor\'ias: I: la inflaci\'on desaparecer\'a antes del cambio de gobierno; II: ocurrir\' una depresi\'on y III: llegar\'a una recesi\'on. Estiman que las probabilidades de que se materealicen las tres teor\'ias son, respectivamente: 0.4, 0.35 y 0.25. As\' mismo, consideran que las probabilidades de que este pa\'is salga del subdesarrollo, si ocurren los eventos son de, 0.90, 0.60 y 0.20 en ese orden. Supongamos que el pa\'is de todos modos no sale del subdesarrollo  cu\'al es la probabilidad de que la inflaci\'on haya desaparecido antes del cambio de gobierno?

\item La probabilidad de que una mujer que da a luz por primera vez tenga un beb\'e con alg\'un s\'indrome o defecto cong\'enito depende de muchos factores; entre otros la edad.  La revista { \em Medical Newsletter} julio de 1999 public\' el siguiente cuadro con estad\'isticas de quienes daban a luz por primera vez

\begin{table}[!ht]
\begin{center}
\begin{tabular}{|c|c|c|c|}
\hline
 & Edad de la Mujer  & Porcentaje de Mujeres & Probabilidad de alg\'un defecto cong\'enito \\\hline
$A_{1}$ & 15 o menos & 3 & 0.05 \\\hline
$A_{2}$ & 16 a 22 & 23 & 0.007 \\\hline
$A_{3}$ & 23 a 29 & 55 & 0.001 \\\hline
$A_{4}$ & 30 a 36 & 12 & 0.04 \\\hline
$A_{5}$ & 37 a 43 & 6 & 0.17 \\\hline
$A_{6}$ & 44 o m\'s & 1 & 0.23 \\\hline
\end{tabular}
\end{center}
\end{table}

De acuerdo con tales datos, si el primer beb\'e naci\'o con alg\'un defecto cong\'enito, cu\'al es la probabilidad de que la edad de la se\~nora oscile entre los 37 y los 43 a\~nos.

\item Un ingeniero fabrica piezas de ajedr\'ez de pl\'astico. Las acabadas se van metiendo en tres grandes cajas antes de pegarles fieltro en la base. En la caja 1 hay $100\%$ de piezas de color blanco, en la 2 hay $50\%$ de cada color, y en la 3 hay $100\%$ de negras. el ingeniero encarga a un empleado que baje la caja de las piezas blancas para pegarles el fieltro, pero como las tres est\a'n muy altas, qui\'en har\'a el trabajo no alacanza a ver su contenido, as\'i que saca una pieza de una caja al azar, que resulta ser de color blanco. Cu\'al es la probabilidad de que sea la caja que busca?


\item Un ni\~no usa calcetines s\'olo de dos colores: azul y negro. Pero no los tiene ordenados por parejas, sino sueltos en dos cajones de su ropero. En el de arriba hay 6 calcetines negors y 2 azules; en el de abajo, 3 negros y 5 azules. No puede encender la luz para ver, porque despertar\'ia a su hermano menor, as\'i que en la oscuridad toma un calcet\'in de cada caj\'on, se los pone, se viste y se va a la escuela.  Cu\'al es la probabilidad de que se haya puesto calcetines del mismo color?

\item La compa\~n\'ia Lear usa 4 empresas para transporte: A, B, C y D. Se sabe que $20\%$ de los embarques se asignan a la empresa A, $250\%$ a la empresa B, $40\%$ a C y $15\%$ a la D. Los embarques llegan retrasados a sus clientes en $7\%$ si los entrega A, $8\%$ si es B, $5\%$ si es C y $9\%$ si es D. Si sabemos que el embarque de hoy fue entregado con retraso,  cu\'al es la probabilidad de que haya sido entregado por la empresa A?

\item Don gato tiene tres candidatos: Dem\'ostenes, Cucho y Benito Bodoque, para una misi\'on que consiste en vender boletos de una rifa inexistente a unos turistas ingenuos. S\'olo uno de los tres debe hacerlo. Estima que Dem\'ostenes cuenta con $35\%$ de ser elegido, Cucho $45\%$ y Benito $20\%$. Hay una probabilidad de $0.09$ de que los atrape el Sargento Matute si va Dem\'ostenes; de $0.11$ si es Cucho, y $0.07$ si Benito es el escogido. Si Matute lo atrap\'o, cu\'al es la probabilidad de que haya sido elegido Benito como el encargado de vender los boletos?

\item La bella y joven Paola es asediada por dos pretendientes: Fernando y Ricardo. Ella est\' decidida a ser novia de alguno de los dos. Fernando tiene una probabilidad de 0.7 de ser el elegido y Ricardo de 0.3. Si el primero resulta afortunado, hay una probabilidad de 0.40 de que terminen en matrimonio; si fuera el segundo, de 0.30. Si Paola se cas\'i con uno de ellos,  cu\'al es la probabilidad de que haya sido Ricardo?
\end{enumerate}
\end{Ejer}




\subsection{Miscel\'anea de ejercicios}


\subsubsection*{Operaciones entre conjuntos}

\begin{Ejer} Resolver los siguientes ejercicios
\begin{enumerate}

\item Proporcione las leyes de D'Morgan para conjuntos.


\item Se desea realizar una rifa para recaudar fondos para asistir a un evento, la cantidad m\'inima por recolectar es de 15000 y se desea que el valor m\'aximo por boleto sea de 100 pesos. Se requiere utilizar solamente 4 d\'igitos. Cu\'antos (posibles combinaciones) n\'umeros es posible generar.


\item En un torneo del deporte m\'as popular se requieren 21 untos para acceder a la siguiente ronda, son 17 equipos.  Cu\'ales
son los resultados posibles que aseguran la calificaci\'on si dan 3 puntos por victoria, 1 por empate y 0 por derrota?


\item En un auditorio con capacidad para 450 personas se desea hacer un evento con 20 filas y dos cabeceras. Para acceder al evento se requiere que satisfagan 6 de las 10 condiciones solicitadas.  Cu\'antas posibles formas hay de seleccion si hay 3 condiciones indispensables: cu\'antas maneras distintas hay de elegir a los solicitantes? Se piensa aceptar exactamente la misma cantidad de hombres que mujeres. Cu\'antas maneras distintas posibles hay de sentarlas si se quiere que haya:  En cada fila la misma cantidad de hombres que mujeres; de manera alternada una fila de hombres y una de mujeres; siempre en cada fila un hombre seguido de una mujer; el arreglo (H=Hombre, M=Mujer): $HHMMHHMMHH$; el arreglo HMMMHMMMH.


\item En el pr\'oximo proceso electoral se inscriben 3 parejas: HM; HH y MM. Hay 500 votantes de los cuales 350 son mujeres y 150 hombres. Suponiendo que pueden votar por 2 duplas: Cu\'antos votos puede obtener cada pareja? Cu\a'ntos votos m\'inimos requiere una dupla para ganar, si son electas 2 de las 3 parejas?


\item Una emisora universitaria cuenta con m\'usica de 2 generos, de los cuales hay 9 discos de genero 1 y 27 de genero 2. En promedio cada disco tiene 13 canciones: De cu\'antas maneras posibles puede programar la m\'usica?   Si en promedio las canciones duran 3.5 minutos cuantos programas de 1 hora aproximandamente puede organizar?  Si desea hacer una vez a la semana un programa especial de cada genero con duraci\'on de 1.5 horas.  Cu\'antas maneras distintas hay de realizar los programas y cuantos son?.


\item El encargado de sistemas del plantel Casa Libertad piensa generar codigos hexadecimales considerando las \'areas existentes: verde, blanca, amarilla, azul, morada, cafe y naranja. Si hay 10 cub\'iculos en cada una de ellas con dos nodos cada uno. Cu\'antos c\'odigos posibles se pueden generar si el primer elemento debe distingir la zona y el segundo debe incluir el numero de cub\'iculo, y as\'i sucesivamente deben estar considerados todos los elementos descritos.: El area verde tiene m\'as personas laborando, lo mismo que la blanca.  En el \'area azul hay 5 laboratorios con 20 m\'aquinas cada uno.


\item Elaborar c\'odigos con 3 numeros y 3 letras, el c\'odigo no puede comenzar con $0$ o la letra \textit{o}.



\item Un adulto mayor de 50 a\~nos se selecciona al azar en una comunidad, en la cual $9\%$ de quienes rebasan esa edad sufren de diabetes, por lo que se les somete a una prueba simple de nivel de glucosa para detectar o desechar la presencia del padecimiento. Sin embargo, el examen no es totalmente confiable, pues a $3\%$ de las personas que no sufren el mal les se\~nala como positivos, mientras que en $15\%$ de aquellos que s\'i est\'an enfermos, la prueba resulta negativa.

\begin{enumerate}
\item Cu\'al es la probabilidad de que ese individuo tenga realmente diabetes, dado que el resultado marca positivo?
\item Cu\'aal es la probabilidad de que no padezca ese mal si marca negativo?
\end{enumerate}

\item En cierto lugar, $30\%$ de las personas son fumadoras y $70\%$ no fumadoras. Adem\'as se estima que $60\%$ de los fumadores y s\'olo $20\%$ de los no fumadores desarrollan hipertensi\'on. De los fumadores hipertensos, $90\%$ llegan a sufrir problemas cardiacos; de los no hipertensos, s\'olo $15\%$ los manifiestan. En cambio, de los no fumadores hipertensos, $65\%$ llega a padecer malestares cardiacos, y de los no hipertensos, s\'olo $5\%$. Si a un individuo se le diagnostica un malestar cardiaco,  cu\'al es la probabilidad de que sea no fumador hipertenso?

\end{enumerate}
\end{Ejer}


\begin{Ejer} Resolver la siguiente lista de ejercicios:
\begin{enumerate}

\item En una f\'abrica hay m\'aquinas de helado que producen $50\%$ y $50\%$ del total.  La A elabora $5\%$ del helado de baja calidad y la B, $6\%$. Encuentre la probabilidad de que un helado de baja calidad provenga de la m\'aquina A.

\item En un programa familiar de televisi\'on, el anfitri\'on le da al concursante la opci\'n de escoger entre las cajas $C_{1}$ y $C_{2}$, con id\'entica apariencia que supuestamente contiene dinero en billetes. El aspirante tiene que seleccionar una, meter la mano y extraer un billete al azar. S\'olo el anfitri\'on sabe que en la caja $C_{1}$ hay seis billetes de 200 pesos y dos de 500, mientras que en la $C_{2}$ hay cinco de 200 y tres de 500. Cuando el concursante saca el billete de una de las cajas, el maestro de ceremonias estba ditra\'ido saludando a alguien del p\'ublico; en ese momento, el primero le muestra el billete que es de 500 pesos. Cu\'al es la probabilidad de que lo haya sacado de la caja $C_{2}$
\end{enumerate}


\end{Ejer}



\subsubsection*{Operaciones entre intervalos}
\begin{Ejer}
Para los siguientes intervalos calcula $A\cap B$, $A\cap C$, $B\cap C$, $A\cup B$, $B\cup C$, $A-B$, $A-C$, $B-C$, $B-A$, $C-B$, $A-(B\cap C)$, $A-(B\cup C)$, $A-(B\cap C)$, $B-(A\cap C)$, $B-(A\cup C)$, $A-B-C)$, $(A\cap B)^{c}$, $(A\cap C)^{c}$, $(C\cap B)^{c}$, $(C\cup B)^{c}$, $A\cap B \cap C$.
\begin{multicols}{2}
\begin{enumerate}
\item $A= [0,2], B= [1,3], C= [2,4]$

\item $A= [-1,1]. B= [0,2], C= [1,3]$

\item $A= [0,1],B= [0,2],C= [1,2]$

\item $A= [0,1], B= [-1,0], C=[-2,-1]$

\item  $A= [-1,1], C= [-2,2], C=[-3,3]$

\item $A=(-5, 0),B=(0, 5),C=(5, 10)$

\item $A=(0, 1),B=(1.5, 2.5),C=(3, 4)$

\item $A= (-\infty, -1),B=(1, 2),C=(3, \infty)$

\item $A= (0, \frac{1}{2}),B= (\frac{1}{2}, 1), C= (1, 2)$

\item $A= (-3, \frac{3}{4}),B= (\frac{2}{5}, 3), C= (-1, 1)$

\item $A= [0,1]$, $B= [-5,5]$, $C= [2,6]$

\item $A= [-\pi,\pi]$, $B= [-3,-1]$, $C=[10,15]$

\item $A= [-\sqrt{2},\sqrt{2}]$, $B= [-\frac{3}{4},\frac{3}{4}]$, $C=[-10,10]$

\item $A= [-2, 0],B=[0, 3],C= [3, 5]$

\item $A= [1, 4],B=[4, 5],C=[5, 7]$

\end{enumerate}
\end{multicols}
\end{Ejer}


\end{document}